\section*{Abstract}

This thesis introduces the application of Deep Equilibrium Nets (DEQNs) to the Analytic Climate Economy (ACE), aiming to lift the curse of dimensionality present in traditional solution methods. To the best of my knowledge, it is one of the first applications of this solution method to integrated assessment models (IAMs). It is also one of the first to contextualize this method in terms of the reinforcement learning literature. With that, this research is a step toward making it possible to relax restrictive assumptions typically necessary in existing models. In this revised version, with corrected equilibrium conditions and a scaling of the objective functions, the policy network recovers the analytical solution of ACE with the DICE 2016 and DICE 2023 production functions closely. The estimated consumption rate deviates from the analytical rate by less than 0.3 percentage points, the abatement rate by 0.1 to 0.2 percentage points on average, and the implied social cost of carbon deviates from its closed-form value by 0.2 to 0.6\% on average. The results also show how sensitive the method is to the correctness of the equilibrium conditions: errors in the submitted version led the network to a wrong fixed point, which was mistakenly attributed to the normalization of inputs. The thesis suggests exploring more effective algorithms inspired by advances in reinforcement learning.